\section{From mode-frame transport to gauge-covariant envelope dynamics}
\label{sec:effective-field}

\subsection{Multiscale ansatz and retained polarization subspace}

Consider a band-isolated excitation whose local carrier angular frequency is $\omega_0$ (determined per wavepacket
by the linearized carrier operator about the slowly varying background at that location).
Write the substrate displacement as a slowly varying complex envelope multiplied by a local mode frame:
\begin{equation}
  \vecu(x,t)\;\approx\;
  \mathrm{Re}\!\left[
    \e^{-\ii\omega_0 t}\sum_{a=1}^{n}\Psi^a(x,t)\,\mathbf{e}_a(x,t)
  \right],
  \label{eq:multiscale-ansatz}
\end{equation}
where $\Psi(x,t)\in\C^n$ varies slowly compared to $\omega_0^{-1}$ and
$\{\mathbf{e}_a(x,t)\}_{a=1..n}$ is an orthonormal frame spanning the retained polarization subspace.
The frame is not unique: $\mathbf{e}_a\mapsto \sum_b \mathbf{e}_b\,U_{ba}(x)$ with $U(x)\in U(n)$ represents the same
physical subspace.

Projecting the full substrate dynamics onto the retained subspace produces, generically, a gauge-covariant derivative
acting on $\Psi$:
\begin{equation}
  \mathcal{D}_\mu \Psi
  =
  \left(\partial_\mu + \ii\,\mathcal{A}_\mu\right)\Psi,
  \label{eq:covD}
\end{equation}
where $\mathcal{A}_\mu$ is the Berry (rank-1) or Wilczek--Zee (non-Abelian) connection defined in \Cref{sec:geophase}.
The gauge transformation $\mathbf{E}\mapsto \mathbf{E}U$ induces $\Psi\mapsto U^\dagger\Psi$ and the standard
connection transformation law, leaving $\mathcal{D}_\mu\Psi$ covariant.

\subsection{Effective action for the slow sector}

Under the multiscale separation (single-band or near-degenerate subspace) and smoothness of the polarization bundle,
the coarse-grained slow dynamics admits a gauge-invariant effective action built from $\Psi$ and the curvature
$\mathcal{F}_{\mu\nu}$:
\begin{equation}
  S_{\text{eff}}[\Psi,\mathcal{A}]
  \;=\;\int \dd^4x\;
  \Big(
    \Psi^\dagger \mathcal{K}(\mathcal{D}) \Psi
    \;-\;\kappa\,\mathrm{tr}(\mathcal{F}_{\mu\nu}\mathcal{F}^{\mu\nu})
    \;+\;\cdots
  \Big),
  \label{eq:Seff}
\end{equation}
where $\mathcal{K}(\mathcal{D})$ is an operator determined by the linearized carrier band (e.g.\ a Klein--Gordon-type or
Dirac-type envelope operator), $\kappa$ is an emergent stiffness for curvature fluctuations, and ``$\cdots$'' denotes
higher-order and symmetry-allowed terms.
In the Abelian case ($n=1$) the curvature term reduces to $F_{\mu\nu}F^{\mu\nu}$ with $F_{\mu\nu}=f_{\mu\nu}$.
For comparison with electromagnetic conventions one may introduce a dimensionful potential
$A^{\mathrm{eff}}_\mu := \eta\,a_\mu$ (and analogously in the non-Abelian case),
so that $F^{\mathrm{eff}}_{\mu\nu}=\partial_\mu A^{\mathrm{eff}}_\nu-\partial_\nu A^{\mathrm{eff}}_\mu=\eta\,f_{\mu\nu}$.
The conversion factor $\eta$ fixes units and would have to be determined by matching to experiment; it is not introduced
as a fundamental coupling in the microscopic substrate equations.

\subsection{\texorpdfstring{Triplet envelope sector and $U(3)\simeq U(1)\times SU(3)$ decomposition}{Triplet envelope sector and U(3) ~= U(1) x SU(3) decomposition}}
\label{subsec:triplet-u3}

If the retained band-isolated sector is a three-dimensional near-degenerate subspace (as suggested by an axis-aligned
triplet on a cubic lattice), then $n=3$ and the emergent connection is $\mathcal{A}_\mu\in\mathfrak{u}(3)$.
Decomposing into trace and traceless parts,
\begin{equation}
  \mathcal{A}_\mu
  = \underbrace{\frac{1}{3}\mathrm{tr}(\mathcal{A}_\mu)\,I_3}_{U(1)\ \text{part}}
  + \underbrace{\sum_{a=1}^{8}A^a_\mu\,T^a}_{SU(3)\ \text{part}},
\end{equation}
suggests a direct route to a ``QCD-like'' internal gauge structure in the slow-sector EFT:
the $U(1)$ part behaves electromagnetism-like, while the traceless part carries non-Abelian self-couplings through
the curvature commutator term in $\mathcal{F}_{\mu\nu}$.
The pairing of an Abelian factor with a non-Abelian one, and the normalization relating them, is the familiar
structure of unified gauge theory \citep{GeorgiGlashow1974,BaezHuerta2010GUTAlgebra}.
Crucially, for nonzero inter-axis mixing, statements about ``a Berry phase per axis'' are not gauge-invariant; the
invariant content is the holonomy/curvature of the transported subspace (Wilczek--Zee).
In that sense, the triplet sector is the natural effective setting in which a baryon-like confined mode should first be
searched numerically on the cubic substrate.

\paragraph{EM as the symmetric (trace) direction.}
The $U(1)$/trace component of $\Psi=(\psi_x,\psi_y,\psi_z)^\top$ is the projection onto the
symmetric triplet direction $(1,1,1)/\sqrt{3}$:
\begin{equation}
  \Psi_{\mathrm{EM}} \;=\; \frac{1}{\sqrt{3}}(\psi_x+\psi_y+\psi_z)\,\frac{(1,1,1)^\top}{\sqrt{3}}.
  \label{eq:em-trace}
\end{equation}
A pure-EM (charge-carrying) excitation therefore displaces all three lateral axis components
\emph{in phase} with equal weight; the traceless complement ($SU(3)$) encodes differential
phase and amplitude relationships \emph{between} the three components.
As a consequence, a single-component displacement (e.g.\ $\psi_x\neq0$, $\psi_y=\psi_z=0$)
projects as $\tfrac{1}{3}$ into the $U(1)$ trace and $\tfrac{2}{3}$ into the $SU(3)$-traceless
sector; it is therefore not a pure-EM excitation.

\paragraph{Topological structure of the two sectors.}
The $U(1)$ and $SU(3)$ factors carry sharply different topology, which is a mathematical fact about
the target spaces of the respective $\Psi$-sectors
\citep{Kibble1976Topology,VilenkinShellard1994,MantonSutcliffe2004}:
\begin{itemize}
  \item $\pi_1(U(1)) = \mathbb{Z}$: the $U(1)$/EM sector admits \emph{line vortex defects} (codimension~2
    in the 3D slice; a 2D worldsheet in the 4D ambient).
    A $U(1)$ vortex is a configuration in which the carrier phase winds by $2\pi$ around the
    vortex core \citep{NielsenOlesen1973,HindmarshKibble1995CosmicStrings}.
  \item $\pi_1(SU(3)) = 0$ (simply connected): the $SU(3)$/color sector admits \emph{no} vortices.
  \item $\pi_3(SU(3)) = \mathbb{Z}$: the $SU(3)$ sector admits \emph{textures} (Skyrmion/instanton-type,
    codimension~3; a localized lump / worldline in 4D) \citep{Skyrme1962,AtiyahManton1989}.
\end{itemize}
The two emergent charge sectors therefore support topologically \emph{different} classes of
localized objects: codim-2 vortex worldsheets in the EM/$U(1)$ sector versus codim-3 texture
worldlines in the color/$SU(3)$ sector.
Field theory provides explicit precedents both for the coexistence of these sectors and for the subtlety of
their stability: a gauged $U(1)$ vortex can carry internal non-Abelian orientational moduli (semilocal and
non-Abelian vortices), and such defects can be stable for dynamical rather than purely topological reasons---e.g.\
semilocal strings are stable on a simply connected vacuum below a critical scalar-to-vector mass ratio
\citep{VachaspatiAchucarro1991Semilocal,AchucarroVachaspati2000,AuzziBolognesiEvslinKonishiYung2003,Eto2021OrientationalModuli}.
This is a structural observation about the target space of the $U(3)$ field; which of these
topological objects the substrate dynamics actually stabilizes, and the identification of specific
objects with physical particles, are targets of the numerical program rather than conclusions of
the present paper.
(Open investigations under \texttt{OPEN\_PROBLEMS.md} D4 and \texttt{EXPERIMENT.md}.)

Geometric phases are universal in wave systems \citep{Pancharatnam1956,Berry1984,WilczekZee1984,Bliokh2015SpinOrbit,Cisowski2022}.
What is specific here is the identification of the physically relevant gauge structure with the
connection induced by adiabatic transport of a band-isolated wavepacket's eigenbundle on the substrate.

\subsection{Dirac-envelope interpretation and the ``split''}

In standard quantum field theory, the Dirac field $\psi$ simultaneously encodes (i) the internal particle labels
(spin, particle/antiparticle) and (ii) the local phase/gauge coupling through the covariant derivative
$D_\mu=\partial_\mu+\ii(q/\hbar)A_\mu$.
In the present theory, this mixture is reinterpreted as follows:
\begin{itemize}
\item the \emph{particle sector} corresponds to localized solitonic modes with additional internal structure
(topology/geometry and its associated labels);
\item the \emph{phase/gauge sector} corresponds to the Berry/Wilczek--Zee connection induced by adiabatic transport
of the band-isolated wavepacket's retained polarization subspace (its local eigenbundle).
\end{itemize}
An effective Dirac-type envelope equation is then understood as a compact coupling of these two sectors.
In particular, a ``charge''-like coupling parameter is not introduced as a fundamental constant of a new field;
it appears as an effective conversion factor between the normalization conventions of the envelope field and the
geometric connection extracted from the substrate.
That a single parameter can set the relative normalization between Abelian and non-Abelian gauge sectors is itself
familiar---e.g.\ the Abelian kinetic-mixing parameter \citep{Holdom1986KineticMixing,Fabbrichesi2021DarkPhoton} and
the single charge-embedding parameter that fixes the Abelian-to-non-Abelian coupling ratio in extended gauge groups
\citep{Hati2017LeftRight3331}.
